\chapter{Backpropagation Through the Actual Architecture}
\label{chap:ch14}

\section{The Computational Graph: Forward Pass and Autograd Engine}
When PyTorch executes the forward pass of \texttt{ECRformerAdaptedModel}, it dynamically constructs a directed acyclic graph (DAG) of tensor operations in memory. 

Each tensor created by a mathematical operation retains a reference to its creator function via the \texttt{grad\_fn} attribute (e.g., \texttt{<AddBackward0>}, \texttt{<ConvolutionBackward0>}, \texttt{<MulBackward0>}). When \texttt{loss.backward()} is invoked on the final scalar loss, PyTorch traverses this DAG in reverse topological order, applying the multivariable Chain Rule to compute:
\begin{equation}
\mathbf{g}_\theta = \frac{\partial \mathcal{L}}{\partial \theta} \quad \text{for every trainable parameter } \theta
\end{equation}

Figure~\ref{fig:backprop_graph} illustrates the forward flow of activations and the reverse propagation of gradients through our verified network architecture.

\begin{figure}[htbp]
\centering
\begin{tikzpicture}[
    font=\footnotesize\sffamily,
    fwd/.style={draw=blue!80, thick, rounded corners=3pt, fill=blue!10, minimum width=2.4cm, minimum height=0.7cm, align=center},
    loss/.style={draw=orange!80, thick, rounded corners=3pt, fill=orange!15, minimum width=2.4cm, minimum height=0.7cm, align=center},
    fwdarr/.style={draw=blue!80, thick, -Stealth},
    backarr/.style={draw=red!80, thick, dashed, -Stealth}
]

% Nodes left to right
\node[fwd] (stem) at (0, 2.5) {\textbf{Stem}\\SAR $\oplus$ Opt};
\node[fwd] (enc) at (3.0, 2.5) {\textbf{Encoder}\\Stages 0--2};
\node[fwd] (dec) at (6.0, 2.5) {\textbf{Decoder}\\Stages 0--2};
\node[fwd] (head) at (9.0, 2.5) {\textbf{Output Head}\\$\mathbf{P}_{\text{dir}}$};
\node[fwd] (blend) at (11.8, 2.5) {\textbf{Soft Blend}\\$\hat{\mathbf{Y}}$};

% Losses
\node[loss] (sdfl_enc) at (3.0, 0.5) {$\mathcal{L}_{\text{enc\_aux}}$\\(\texttt{down\_proj})};
\node[loss] (sdfl_dec) at (6.0, 0.5) {$\mathcal{L}_{\text{dec\_aux}}$\\(\texttt{up\_proj})};
\node[loss] (main_loss) at (11.8, 0.5) {$\mathcal{L}_{\text{composite}}$\\($L_1$ + VGG + GAN)};

% Forward arrows (Blue)
\draw[fwdarr] (stem) -- (enc);
\draw[fwdarr] (enc) -- (dec);
\draw[fwdarr] (dec) -- (head);
\draw[fwdarr] (head) -- (blend);
\draw[fwdarr] (blend) -- (main_loss);
\draw[fwdarr] (enc) -- (sdfl_enc);
\draw[fwdarr] (dec) -- (sdfl_dec);

% Backward gradient arrows (Red dashed)
\draw[backarr] (main_loss) -- (blend) node[midway, right=2pt, font=\scriptsize, text=red!80!black] {$\frac{\partial \mathcal{L}}{\partial \hat{\mathbf{Y}}}$};
\draw[backarr] (blend) -- (head) node[midway, below=2pt, font=\scriptsize, text=red!80!black] {$\times \mathbf{M}$};
\draw[backarr] (head) -- (dec) node[midway, below=2pt, font=\scriptsize, text=red!80!black] {$\nabla_{\mathbf{W}_{\text{head}}}$};
\draw[backarr] (dec) -- (enc) node[midway, below=2pt, font=\scriptsize, text=red!80!black] {$\nabla_{\mathbf{W}_{\text{dec}}} + \text{skips}$};
\draw[backarr] (enc) -- (stem) node[midway, below=2pt, font=\scriptsize, text=red!80!black] {$\nabla_{\mathbf{W}_{\text{enc}}}$};

% Auxiliary gradient injections
\draw[backarr] (sdfl_dec) -- (dec) node[midway, left=2pt, font=\scriptsize, text=red!80!black] {$+ \nabla \mathcal{L}_{\text{aux}}$};
\draw[backarr] (sdfl_enc) -- (enc) node[midway, left=2pt, font=\scriptsize, text=red!80!black] {$+ \nabla \mathcal{L}_{\text{aux}}$};

% Note on SAR stem gradient scaling
\node[draw=red!70, dashed, fill=red!5, rounded corners=3pt, align=center, font=\scriptsize] at (0, 0.5) {\textbf{$2\times$ Learning Rate}\\applied to Branch 0\\(\texttt{stem.branch.0})};
\draw[backarr] (0, 1.0) -- (stem.south);

\end{tikzpicture}
\caption{Forward activation flow (solid blue) and backward gradient propagation (dashed red) through \texttt{ECRformerAdaptedModel}. The soft-mask blend gate multiplies gradients from the main loss by $\mathbf{M}$, ensuring that gradients to the generator are concentrated on cloud-covered regions. Intermediate auxiliary heads inject direct gradient signals into earlier decoder and encoder stages.}
\label{fig:backprop_graph}
\end{figure}


\section{Gradient Flow Through the Soft-Mask Blend Gate}
One of the most elegant mathematical properties of our architecture occurs at the final output blending layer:
\begin{equation}
\hat{\mathbf{Y}} = (1 - \mathbf{M}) \odot \mathbf{X}_{S2} + \mathbf{M} \odot \mathbf{P}_{\text{dir}}
\end{equation}
When computing the gradient of the composite loss $\mathcal{L}$ with respect to the generator's direct prediction $\mathbf{P}_{\text{dir}}$:
\begin{equation}
\frac{\partial \mathcal{L}}{\partial \mathbf{P}_{\text{dir}}} = \frac{\partial \mathcal{L}}{\partial \hat{\mathbf{Y}}} \odot \frac{\partial \hat{\mathbf{Y}}}{\partial \mathbf{P}_{\text{dir}}} = \frac{\partial \mathcal{L}}{\partial \hat{\mathbf{Y}}} \odot \mathbf{M}
\end{equation}

\begin{conceptbox}{Mathematical Significance: The Zero-Gradient Shield}
Look carefully at the term $\frac{\partial \mathcal{L}}{\partial \mathbf{P}_{\text{dir}}} = \frac{\partial \mathcal{L}}{\partial \hat{\mathbf{Y}}} \odot \mathbf{M}$:
\begin{itemize}
    \item \textbf{Under Cloudy Pixels ($\mathbf{M} \approx 1.0$):} $\frac{\partial \mathcal{L}}{\partial \mathbf{P}_{\text{dir}}} \approx \frac{\partial \mathcal{L}}{\partial \hat{\mathbf{Y}}}$. Full gradient magnitude flows back to the neural network weights, driving learning and reconstruction.
    \item \textbf{Under Clear Pixels ($\mathbf{M} = 0.0$):} $\frac{\partial \mathcal{L}}{\partial \mathbf{P}_{\text{dir}}} = \frac{\partial \mathcal{L}}{\partial \hat{\mathbf{Y}}} \times 0.0 = \mathbf{0.0}$!
\end{itemize}
The generator receives \textbf{zero gradient penalty} for what it predicts under clear skies! The soft mask acts as a mathematical shield: the network is not forced to waste parameter capacity memorizing clear terrain that will simply be preserved from the original optical image anyway!
\end{conceptbox}

\section{Gradient Propagation Across Skip Connections and Auxiliary Heads}
As the gradient travels backward through the decoder, two critical structural bifurcations occur:

\subsection{1. Additive Gradient Injection from Auxiliary Heads}
At each decoder level $k$, the feature map $\mathbf{F}_{\text{dec}, k}$ feeds both into the main decoder path and into an auxiliary projection head (\texttt{up\_proj[k]}):
\begin{equation}
\frac{\partial \mathcal{L}_{\text{total}}}{\partial \mathbf{F}_{\text{dec}, k}} = \frac{\partial \mathcal{L}_{\text{main}}}{\partial \mathbf{F}_{\text{dec}, k}} + w_{\text{aux}} \frac{\partial \mathcal{L}_{\text{aux}, k}}{\partial \mathbf{F}_{\text{dec}, k}}
\end{equation}
The gradients add together directly! This injects strong, direct supervision signals into early decoder and encoder layers, bypassing the vanishing gradient problem.

\subsection{2. Bifurcation at Skip Connections}
When gradients reach a concatenation node with an encoder skip connection:
\begin{equation}
\mathbf{F}_{\text{fused}} = [\mathbf{F}_{\text{upsampled}}, \mathbf{F}_{\text{skip}}]
\end{equation}
The incoming gradient vector splits along the channel axis:
\begin{itemize}
    \item Gradient slice $[0 : C_{\text{up}}]$ flows backward into the upsampling deconvolution path.
    \item Gradient slice $[C_{\text{up}} : 2C_{\text{up}}]$ jumps directly across the skip connection into the encoder!
\end{itemize}
The encoder receives gradients through two parallel highways: directly via skip connections and indirectly through the bottleneck.

\section{The $2\times$ Differential Learning Rate on SAR Stem}
In \texttt{train\_ecrformer\_v3.py} and \texttt{train\_ecrformer\_v4\_gan.py}, we implemented a specialized parameter grouping in the optimizer:
\begin{lstlisting}[language=Python]
sar_params = list(model.module.stem.branch_0.parameters()) + \
             list(model.module.encoder[0].parameters())
base_params = [p for p in model.parameters() if p not in set(sar_params)]

optimizer = torch.optim.AdamW([
    {'params': base_params, 'lr': 1e-4},
    {'params': sar_params, 'lr': 2e-4} # 2x learning rate for SAR features!
], weight_decay=1e-4)
\end{lstlisting}
\begin{itemize}
    \item \textbf{Why this is necessary:} Microwave radar backscatter has a lower signal-to-noise ratio and higher statistical variance (speckle noise) than optical imagery. In standard training, gradients preferentially update the optical branch, leaving the radar branch under-utilized.
    \item Doubling the learning rate ($2\times 10^{-4}$) on \texttt{branch\_0} forces the network to aggressively learn high-contrast SAR edge extractors during the earliest epochs!
\end{itemize}

\section{Summary of Chapter 14}
\begin{enumerate}
    \item Autograd constructs a dynamic computational DAG and computes exact partial derivatives via reverse topological traversal.
    \item The soft-mask blend gate multiplies gradients to the generator by $\mathbf{M}$, shielding the network from wasting capacity on clear pixels.
    \item Auxiliary projection heads inject additive gradient signals directly into intermediate layers.
    \item Skip connections provide direct gradient expressways from the loss to early encoder layers.
    \item A $2\times$ differential learning rate forces aggressive feature extraction on Sentinel-1 SAR channels.
\end{enumerate}

\section{Practice Exercises}
\begin{enumerate}
    \item If a pixel has mask value $\mathbf{M} = 0.20$ and the upstream loss gradient $\frac{\partial \mathcal{L}}{\partial \hat{\mathbf{Y}}} = 5.0$, what is the gradient received by the generator prediction $\frac{\partial \mathcal{L}}{\partial \mathbf{P}_{\text{dir}}}$?
    \item Why does applying auxiliary losses to intermediate encoder/decoder layers prevent vanishing gradients?
    \item Explain why the SAR stem branch (\texttt{branch\_0}) was assigned twice the learning rate of the rest of the network.
\end{enumerate}

\subsection*{Solutions}
\begin{enumerate}
    \item $\frac{\partial \mathcal{L}}{\partial \mathbf{P}_{\text{dir}}} = \frac{\partial \mathcal{L}}{\partial \hat{\mathbf{Y}}} \times \mathbf{M} = 5.0 \times 0.20 = 1.0$.
    \item In very deep networks, gradients traveling through dozens of non-linear layers can shrink exponentially toward zero (vanishing gradients). Intermediate auxiliary heads connect the loss function directly to early layers through only a single $3\times 3$ convolution, injecting unattenuated gradient signals directly into early representations.
    \item SAR data contains complex speckle noise and different physical scattering mechanisms that are harder to map to optical reflectance than optical-to-optical features. A $2\times$ learning rate overcomes this inertia, forcing the network to actively exploit microwave information rather than taking the lazy path of relying solely on optical context.
\end{enumerate}
